\subsection{复数的定义}

多项式 \(X^{2} + 1\) 在 R 中没有根，因为 \(x^{2} + 1 \geqslant 1\) 对一切 \(x \in R\) 成立；因此它在 R 上不可约。{[}这是对任何有序域都成立的类似结果的一个特例。{]}

定义 1. 域 \(\mathbf{R}[\mathbf{X}] / (\mathbf{X}^2 + 1)\) 称为复数域，记作 \(\mathbf{C}\) 。\(\mathbf{X}\) 在 \(\mathbf{C}\) 中的典范像记作 \(i\) ，于是 \(\mathbf{C}\) 是由域 \(\mathbf{R}\) 通过代数添加根 \(i\) （多项式 \(\mathbf{X}^2 + 1\) 的根）而得到的。\(\mathbf{C}\) 的元素称为复数。

从代数的观点看，域 \(\mathbf{C}\) 的重要性基于下面的基本定理：

定理 1. （达朗贝尔–高斯）复数域 \(\mathbf{C}\) 是代数闭的。

为证明定理，只需证明：(i) R 中每个 \(\geqslant0\) 的元素都有平方根，并且 \((\mathrm{ii})\) 每个系数在 R 中的奇次数多项式在 R 中至少有一个根。第一个断言已在第四章 § 3 第 3 号中证明。至于第二个，若 \(f(\mathbf{X}) = a_{0}\mathbf{X}^{n} + a_{1}\mathbf{X}^{n-1} + \cdots + a_{n}\) 是一个次数 \(n(a_{0} \neq 0)\) 为奇数的实系数多项式，则可写 \(f(x) = a_{0}x^{n}g(x)\) （对 \(x \neq 0\) ），其中

\[
g (x) = \mathrm{I} + \frac {a _ {1}}{a _ {0} x} + \dots + \frac {a _ {n}}{a _ {0} x ^ {n}}
\]

它趋于 \(+1\) （当 \(x\) 趋于 \(+\infty\) 或 \(-\infty\) 时）。因此存在一个数 \(a > 0\) ，使得 \(f(a)\) 与 \(a_0\) 同号，而 \(f(-a)\) 与 \(-a_0\) 同号；
于是由波尔查诺定理（第四章，§ 6，第 1 号，定理 2），\(f\) 在 \([-a, a]\) 中至少有一个根。

注。1) 定理 1 可以不引用有序域理论，而利用域 C 的拓扑性质（将在下面（第 2 号）定义）来证明；见 § 2 习题 2，并见本丛书专论代数拓扑的部分，在那里达朗贝尔–高斯定理将作为关于映射的度的结果的一个推论而出现。

\begin{enumerate}
\def\labelenumi{\arabic{enumi})}
\setcounter{enumi}{1}
\tightlist
\item
  由于 C 在 R 上的次数为 2，由此推出：在同构意义下，C 是 R 本身之外 R 的唯一的代数扩张，并且不存在包含在 C 中且包含 R 的、异于 R 与 C 的域。
\end{enumerate}

我们知道 \(\mathbf{R}\) 可以等同于 \(\mathbf{C}\) 的一个子域，并且每个 \(z \in \mathbf{C}\) 都可以唯一地写成形式 \(x + iy\) ，其中 \(x\) 与 \(y\) 是实数；\(x\) 称为 \(z\) 的实部，记作 \(\Re(z)\) ；\(y\) 称为 \(z\) 的虚部，记作 \(\Im(z)\) 。形如 \(iy\) （ \(y\) 为实数）的复数称为纯虚数。关系 \(x + iy = o\) （ \(x, y\) 为实数）等价于 \(x = o\) 且 \(y = o\) 。

由于 \(i^2 = -1\) ，\(\mathbf{C}\) 的元素（当以其实部与虚部给出时）满足下面的运算规则：

\begin{enumerate}
\def\labelenumi{(\arabic{enumi})}
\tightlist
\item
\end{enumerate}

\[
(x + i y) + (x ^ {\prime} + i y ^ {\prime}) = (x + x ^ {\prime}) + i (y + y ^ {\prime})\tag{2}
\]

\[
(x + i y) (x ^ {\prime} + i y ^ {\prime}) = (x x ^ {\prime} - y y ^ {\prime}) + i (x y ^ {\prime} + x ^ {\prime} y).
\]

特别地，\((x + iy)(x - iy) = x^{2} + y^{2}\in \mathbf{R}\) ，因此，若 \(x + iy\neq 0\) ，则

\[
\frac {1}{x + i y} = \frac {x}{x ^ {2} + y ^ {2}} - i \frac {y}{x ^ {2} + y ^ {2}}.\tag{3}
\]

多项式 \(X^{2}+1\) 在 C 中的第二个根是 -i；因而 C 中除恒等映射外使所有实数不变的自同构，只有把 \(z=x+iy\) 映为 x-iy 的那一个；后者记作 \(\bar{z}\) ，并且（与一般定义相一致）称为与 z 共轭的复数。我们有

\[
\mathcal {R} (z) = \frac {\mathrm{I}}{2} (z + \bar {z}) \quad \text { and } \quad \mathfrak {J} (z) = \frac {\mathrm{I}}{2 i} (z - \bar {z}).
\]

由于这个自同构，若 \(f(z)\) 是一个实系数多项式，则 \(f(\overline{z}) = \overline{f(z)}\) 对一切 \(z \in \mathbf{C}\) 成立。

实数 \(z\bar{z} = x^2 + y^2\) 称为 \(z\) 的代数范数，或在无混淆危险时简称范数；它是一个 \(\geqslant 0\) 的实数，仅当 \(z = 0\) 时才为零。实数 \(\sqrt{z\bar{z}} = \sqrt{x^2 + y^2} \geqslant 0\) 当 \(z\) 为实数时就是 \(z\) 的绝对值，我们仍称它为 \(z\) 的绝对值，并记作 \(|z|\) ，当 \(z\) 是任意复数时。关系 \(|z| = 0\) 等价于 \(z = 0\) 。若 \(z\) 与 \(z'\) 是两个复数，则 \(zz'\) 的共轭是 \(\overline{zz'}\) ，因此 \(|zz'|^2 = zz' \overline{zz'} = |z|^2 |z'|^2\) ，从而 \(|zz'| = |z|. |z'|\) ：乘积的绝对值等于诸因子的绝对值之积。特别地，若 \(z \neq 0\) 且 \(z' = 1/z\) ，则有 \(|1/z| = 1/|z|\) 。

最后，对一切复数 \(z, z'\) ，我们有三角不等式

\[
\left| z + z ^ {\prime} \right| \leqslant | z | + \left| z ^ {\prime} \right|.\tag{4}
\]

